% Inicio del documento
\documentclass{article}

% Paquetes
\usepackage{geometry} % Sirve para ordenar el formato de hoja
\usepackage{tikz} % Sirve para dibujar
\usepackage[dvipsnames]{xcolor} % Sirver para poner colores
\usepackage{graphicx} % Sirve para poner imagenes

% Uso de pacquetes
\geometry{a4paper, margin=0.5cm, top=0.5cm} % Ordenar los bordes
\setlength{\parindent}{0pt} % Quitar la sangria del primer bloque
\setlength{\fboxsep}{0pt} % Te permite agregar padding a la caja

% Estructura
\begin{document}
  % Dibujo de la linea 
  \begin{tikzpicture}[remember picture, overlay] % Guarda medidas y sobretodo
  \fill[Green!20] (current page.north west) rectangle ([xshift=7cm]current page.south west); % Lo primero define el color y el degradado y deespues que va a cubrir
  \draw[line width=1pt, color=black] ([xshift=7cm]current page.north west) -- ([xshift=7cm]current page.south west); % Desde las esquinas poner un punto a cierta distancia
  \end{tikzpicture}

  \begin{minipage}[t]{0.3\textwidth}
    \begin{center}
      \vspace{10pt}
      \begin{tikzpicture}
        \clip (0,0) circle (3cm); % Crea un circulo de radio 3
        \node at (0,0) {\includegraphics[width=6cm]{~/Imágenes/4.jpg}}; % Usa la imagen con diametro 6 
        \draw[line width=4pt, color=black] (0,0) circle (3cm); % Crea un borde de mismo tamaño q el circulo pero de color negro
      \end{tikzpicture}
      \par \vspace{25pt}
      {\huge \textbf{Sebastián Flores}} \\[10pt] {\LARGE Apoyo inventario}
    \end{center}
  \end{minipage}
  \hfill
  \begin{minipage}[t]{0.6\textwidth}
    \vspace{20pt}
    \section*{Perfil}
    Este es el contenido de la barra derecha...
  \end{minipage}

\end{document}
